% supplement_verification.tex
% Supplemental Material: Mathematical Verification of the tau Framework Across Scales
% Compile with: pdflatex supplement_verification.tex

\documentclass[prl,twocolumn,superscriptaddress,longbibliography,floatfix]{revtex4-2}

\usepackage{amsmath}
\usepackage{amssymb}
\usepackage{amsfonts}
\usepackage{amsthm}
\usepackage{bm}
\usepackage{hyperref}
\usepackage{booktabs}
\usepackage{graphicx}

\newtheorem{theorem}{Theorem}
\newtheorem{corollary}[theorem]{Corollary}
\newtheorem{lemma}[theorem]{Lemma}
\newtheorem*{remark}{Remark}

% Convenience macros (consistent with main papers)
\newcommand{\kB}{k_{\mathrm{B}}}
\newcommand{\Tr}{\mathrm{Tr}}
\newcommand{\id}{\mathrm{id}}
\newcommand{\Petz}{\mathcal{R}}
\newcommand{\chan}{\mathcal{N}}
\newcommand{\hilb}{\mathcal{H}}
\newcommand{\code}{\mathcal{C}}
\newcommand{\KL}{D}
\newcommand{\ket}[1]{|#1\rangle}
\newcommand{\bra}[1]{\langle#1|}
\newcommand{\braket}[2]{\langle#1|#2\rangle}
\newcommand{\op}[2]{|#1\rangle\langle#2|}
\newcommand{\abs}[1]{\left|#1\right|}
\newcommand{\norm}[1]{\left\|#1\right\|}
\newcommand{\tauasym}{\tau}
\newcommand{\BC}{\mathrm{BC}}

\begin{document}

\title{Supplemental Material: Mathematical Verification of the
$\tauasym$ Framework Across Scales}

\author{Sheng-Kai Huang}
\email{akai@fawstudio.com}
\affiliation{Independent Researcher}

\date{\today}

\begin{abstract}
We verify the mathematical consistency of the temporal asymmetry
parameter $\tauasym = 1-F$ across four layers:
(1)~quantum channels, where the JRSWW bound is a proven theorem;
(2)~gravitational fields, where three independent routes motivate
a conjecture;
(3)~classical systems, where the bound reduces to Bayesian
retrodiction; and
(4)~stochastic thermodynamics, where the Crooks fluctuation theorem
provides the bridge.
Each layer is assigned an honest status label.
\end{abstract}

\maketitle

%=======================================================================
\section{Layer 1: Quantum Channels}
\label{sec:layer1}
%=======================================================================

\noindent\textbf{Status: [THEOREM] --- mathematically proven.}

\subsection{The JRSWW bound}

For a quantum channel $\chan$ (CPTP map) and states $\rho, \sigma$
with $\rho \ll \sigma$, the quantum relative entropy (QRE)
satisfies the data processing inequality (DPI):
\begin{equation}
D(\rho\|\sigma) \geq D(\chan(\rho)\|\chan(\sigma)).
\label{eq:DPI}
\end{equation}
Define the QRE drop:
\begin{equation}
\Sigma \;\equiv\; D(\rho\|\sigma) - D(\chan(\rho)\|\chan(\sigma))
\;\geq\; 0.
\label{eq:Sigma_def}
\end{equation}
Junge, Renner, Sutter, Wilde, and Winter~\cite{JRSWW2018}
proved the strengthened DPI: there exists a rotated Petz recovery
map $\tilde{\Petz}$ such that
\begin{equation}
F\bigl(\rho,\;\tilde{\Petz}\circ\chan(\rho)\bigr)
\;\geq\; \exp(-\Sigma/2),
\label{eq:JRSWW}
\end{equation}
where $F(\rho,\omega) = \norm{\sqrt{\rho}\sqrt{\omega}}_1$ is
the Uhlmann fidelity~\cite{Uhlmann1976}. An earlier version with
$F^2$ on the left was obtained by Fawzi and
Renner~\cite{FawziRenner2015}.

The temporal asymmetry parameter is defined as:
\begin{equation}
\tauasym \;\equiv\; 1 - F(\rho,\;\Petz\circ\chan(\rho)),
\label{eq:tau_def}
\end{equation}
so the JRSWW bound becomes the $\tauasym$ bound:
\begin{equation}
\tauasym \;\leq\; 1 - \exp(-\Sigma/2).
\label{eq:tau_bound}
\end{equation}

\subsection{The equivalence chain}

The following conditions are equivalent for a channel $\chan$
with reference state $\sigma$~\cite{Petz1988,FawziRenner2015,HJPW2004}:
\begin{equation}
\tauasym = 0
\;\Longleftrightarrow\;
\Sigma = 0
\;\Longleftrightarrow\;
I(A;E|B)_\omega = 0
\;\Longleftrightarrow\;
\text{QMC},
\label{eq:equiv_chain}
\end{equation}
where $\omega_{ABE}$ is the purification of the channel output,
$I(A;E|B)$ is the quantum conditional mutual information
(QCMI), and QMC denotes the quantum Markov chain condition
$\omega_{ABE} = \Petz_{B\to BE}(\omega_{AB})$.

The proof assembles three known results:
\begin{itemize}
\item $\Sigma = 0 \Leftrightarrow$ exact Petz recovery
(Petz~\cite{Petz1988}).
\item $\Sigma = 0 \Leftrightarrow I(A;E|B)=0$
(Hayden--Jozsa--Petz--Winter~\cite{HJPW2004}).
\item $I(A;E|B) = 0 \Leftrightarrow$ approximate recovery with
fidelity controlled by QCMI
(Fawzi--Renner~\cite{FawziRenner2015}).
\end{itemize}

\subsection{Numerical verification}

We numerically tested the $\tauasym$ bound~\eqref{eq:tau_bound}
on three standard quantum channels: amplitude damping
($\gamma \in [0.01, 0.99]$, step 0.01), dephasing
($p \in [0.01, 0.99]$), and depolarizing
($p \in [0.01, 0.99]$). For each of 99 noise
parameters, 20 random qubit states were generated (Haar-uniform),
yielding $3 \times 99 \times 20 = 5{,}940$ test cases.

\emph{Result.}---Zero violations of Eq.~\eqref{eq:tau_bound}
were found. In all 5{,}940 cases,
$\tauasym_{\text{measured}} \leq 1 - \exp(-\Sigma/2)$.

\emph{Dephasing anomaly.}---For the dephasing channel with
maximally mixed reference state $\sigma = I/2$, the identity map
achieves strictly higher fidelity than the standard Petz map.
This is because the Petz map with $\sigma = I/2$ reduces to
$\Petz = \chan^\dagger \circ (\cdot) = \chan \circ (\cdot)$
(the dephasing channel is self-adjoint), so applying the Petz map
means applying the dephasing channel a second time. The identity map,
which does nothing, trivially outperforms double-dephasing.
This does not violate the JRSWW bound, which guarantees only
that the \emph{rotated} Petz map is near-optimal.

%=======================================================================
\section{Layer 2: Gravitational Channels}
\label{sec:layer2}
%=======================================================================

\noindent\textbf{Status: [MOTIVATED CONJECTURE] --- well-supported
but not proven.}

\subsection{The identification}

Paper~2 of this series~\cite{Paper2} proposes the gravitational
entropy production:
\begin{equation}
\Sigma_{\text{grav}} = -\ln(-g_{00})
= \frac{r_s}{r}
\quad\text{(exponential metric)},
\label{eq:Sigma_grav}
\end{equation}
where $r_s = 2GM/c^2$ is the Schwarzschild radius.
Three independent derivations converge on this identification.

\subsection{Route 1: Modular flow}

Dorau and Much~\cite{DorauMuch2025} showed that modular flow in
algebraic QFT derives the Einstein field equations from quantum
relative entropy. Their construction defines a local QRE density
whose integral over a causal diamond gives the total entropy
production. For a static spherically symmetric spacetime,
the fractional QRE loss between the diamond at infinity and one
at radius $r$ evaluates to:
\begin{equation}
\frac{D_{\text{in}} - D_{\text{out}}}{D_{\text{in}}}
= \frac{r_s}{r}.
\label{eq:modular}
\end{equation}
This is exact, not an approximation.

\subsection{Route 2: Gravitational Landauer principle}

Herrera~\cite{Herrera2020} derived the minimum energy cost of
erasing one bit of information at radius $r$ in a gravitational
field, accounting for the Tolman redshift factor:
\begin{equation}
W_{\text{erase}}(r) = \kB T_{\text{local}} \ln 2
= \frac{\kB T_\infty}{\sqrt{-g_{00}}} \ln 2.
\label{eq:landauer_grav}
\end{equation}
The entropy production per erasure, measured from infinity, is:
\begin{equation}
\Sigma_{\text{Landauer}} = -\ln(-g_{00}),
\label{eq:landauer_sigma}
\end{equation}
matching Eq.~\eqref{eq:Sigma_grav}.

\subsection{Route 3: Bosonic loss channel}

Ahmadi and Fuentes~\cite{AhmadiFuentes2014} showed that
a bosonic field mode propagating in a static spacetime
undergoes a pure-loss bosonic channel with transmissivity
$\eta = -g_{00}$. For a pure-loss channel, the QRE drop
between coherent state inputs is~\cite{Wilde2017}:
\begin{equation}
\Sigma_{\text{bosonic}} = -\ln\eta = -\ln(-g_{00}),
\label{eq:bosonic_sigma}
\end{equation}
again matching Eq.~\eqref{eq:Sigma_grav}.

\subsection{Critical caveats}

\emph{No first-principles CPTP map.}---The JRSWW
bound~\eqref{eq:JRSWW} requires a CPTP map $\chan$ satisfying
$\Delta D = \Sigma$. No single, first-principles CPTP map has been
proven to give $\Delta D = -\ln(-g_{00})$ exactly. Each route above
extracts $\Sigma_{\text{grav}}$ from a different physical
construction; none provides a universal gravitational channel.

\emph{Dimensional obstruction.}---Finite-dimensional (qubit)
channels satisfy $\Delta D \leq \ln d$, where $d$ is the Hilbert
space dimension. Since $\Sigma_{\text{grav}} = r_s/r$ is
unbounded as $r \to r_s$, the gravitational channel must be
infinite-dimensional (bosonic). This is consistent with Route~3
but means qubit numerics cannot probe the strong-field regime.

\emph{Saturation is impossible for $\Sigma > 0$.}---The
Golden-Thompson inequality~\cite{GoldenThompson}
$\Tr(e^A e^B) \geq \Tr(e^{A+B})$ makes the JRSWW bound strictly
non-saturating whenever $[\rho, \sigma] \neq 0$ and $\Sigma > 0$.
The exponential metric $g_{00} = -e^{-r_s/r}$ cannot be derived
from bound saturation alone.

\emph{What survives.}---The identification
\begin{equation}
\frac{c_{\text{eff}}}{c} = \exp(-\Sigma_{\text{grav}}/2)
= F_{\text{bound}}
\label{eq:ceff}
\end{equation}
is a mathematical identity following from the definition
$\Sigma_{\text{grav}} = 2|\Phi|/c^2$ and the definition of the
effective speed of light. It does not depend on saturation and
holds regardless of the channel details.

\emph{Alternative route via quantum hypothesis testing.}---Hayashi
and Yamasaki~\cite{HayashiYamasaki2025} define $\Sigma$ directly
from quantum states via
$\Sigma_{\text{grav}} = D(\rho(z_1)\|\rho(z_2))$, where $\rho(z)$
is the field state at spacetime point $z$. This bypasses the need
for a CPTP map entirely, at the cost of requiring well-defined
local states.

%=======================================================================
\section{Layer 3: Classical Limit}
\label{sec:layer3}
%=======================================================================

\noindent\textbf{Status: [THEOREM] --- proven as special case
of Layer~1.}

\subsection{Statement}

\begin{theorem}[Classical $\tauasym$ Bound]
\label{thm:classical}
For any classical channel (stochastic matrix) $P$ with prior
distribution $p$ and posterior $q = Pp$, the temporal asymmetry
parameter satisfies:
\begin{equation}
\tauasym_{\text{cl}} \;\leq\;
1 - \exp\!\bigl(-\Delta D_{\mathrm{KL}}/2\bigr),
\label{eq:tau_classical}
\end{equation}
where $\Delta D_{\mathrm{KL}} = D_{\mathrm{KL}}(p\|s)
- D_{\mathrm{KL}}(Pp\|Ps)$ is the KL divergence drop,
$s$ is the reference distribution, and $\tauasym_{\text{cl}}$
is computed using the Bayesian retrodiction map.
\end{theorem}

\subsection{Proof: reduction from the quantum case}

The proof consists of showing that each quantum object in the JRSWW
bound reduces to its classical counterpart when all states commute.

\emph{Step 1: Petz map $\to$ Bayes' rule.}---For commuting
states (diagonal density matrices), the Petz recovery
map~\cite{Petz1988,Parzygnat2023}
\begin{equation}
\Petz_{\sigma,\chan}(\omega) = \sigma^{1/2}\,
\chan^\dagger\!\bigl(\chan(\sigma)^{-1/2}\,\omega\,
\chan(\sigma)^{-1/2}\bigr)\,\sigma^{1/2}
\label{eq:petz_map}
\end{equation}
reduces to the Bayesian inverse. Explicitly, with
$\sigma = \mathrm{diag}(s_1,\ldots,s_n)$ and stochastic matrix
$P_{j|i}$:
\begin{equation}
\Petz(P,s)_{i|j}
= \frac{s_i\, P_{j|i}}{(Ps)_j}
= P(i|j),
\label{eq:bayes}
\end{equation}
which is Bayes' rule.

\emph{Step 2: QRE $\to$ KL divergence.}---The quantum relative
entropy $D(\rho\|\sigma)
= \Tr[\rho(\ln\rho - \ln\sigma)]$~\cite{Umegaki1962} reduces to
the Kullback--Leibler divergence
$D_{\mathrm{KL}}(p\|s) = \sum_i p_i \ln(p_i/s_i)$
when $\rho = \mathrm{diag}(p)$ and $\sigma = \mathrm{diag}(s)$.

\emph{Step 3: Uhlmann fidelity $\to$ Bhattacharyya
coefficient.}---For probability distributions $p$ and $q$,
the Uhlmann fidelity becomes~\cite{Uhlmann1976}:
\begin{equation}
F(p,q) = \sum_i \sqrt{p_i\, q_i} \equiv \BC(p,q),
\label{eq:BC}
\end{equation}
the Bhattacharyya coefficient.

\emph{Step 4: Assembling.}---Substituting Steps 1--3 into the
JRSWW bound~\eqref{eq:JRSWW} yields Eq.~\eqref{eq:tau_classical}
directly. \hfill$\square$

\subsection{Consequences}

The classical $\tauasym$ bound applies to any system
describable by a stochastic matrix:
\begin{itemize}
\item Thermodynamic processes (heat engines, molecular motors).
\item Neural dynamics (transition matrices between brain states).
\item Chemical reaction networks (master equations).
\end{itemize}
No additional assumptions beyond the existence of a classical
channel and a reference distribution are needed.

%=======================================================================
\section{Layer 4: Stochastic Thermodynamics Bridge}
\label{sec:layer4}
%=======================================================================

\noindent\textbf{Status: [BRIDGE EXISTS] --- mathematically sound
via the Crooks fluctuation theorem.}

\subsection{The problem}

Stochastic thermodynamics uses the Schnakenberg entropy
production~\cite{Schnakenberg1976}
\begin{equation}
\Sigma_{\text{Sch}} = \sum_{i \neq j}
J_{ij}\,\ln\frac{W_{ij}\, p_i}{W_{ji}\, p_j},
\label{eq:Schnakenberg}
\end{equation}
where $W_{ij}$ are transition rates and $J_{ij}$ are probability
currents. This is \emph{not} the same as the JRSWW quantity
$\Sigma_{\text{DPI}}$ from Eq.~\eqref{eq:Sigma_def}. The Schnakenberg
entropy production counts dissipation along edges of a graph; the
JRSWW quantity measures global KL divergence drop under a channel.
Direct identification would be incorrect.

\subsection{Solution: Crooks on path space}

The bridge is provided by the Crooks fluctuation
theorem~\cite{Crooks1999} applied to trajectory space. For a
system driven by protocol $\lambda(t)$, each individual trajectory
$\gamma$ carries a stochastic entropy production:
\begin{equation}
\Sigma_{\text{traj}}(\gamma)
= \ln\frac{P_{\text{fwd}}(\gamma)}{P_{\text{rev}}(\bar{\gamma})},
\label{eq:Sigma_traj}
\end{equation}
where $P_{\text{fwd}}(\gamma)$ is the probability of trajectory
$\gamma$ under the forward protocol and
$P_{\text{rev}}(\bar{\gamma})$ is the probability of the
time-reversed trajectory under the reversed protocol.

The Crooks theorem states~\cite{Crooks1999}:
\begin{equation}
\bigl\langle e^{-\Sigma_{\text{traj}}} \bigr\rangle = 1.
\label{eq:crooks_average}
\end{equation}

Now treat the forward and reverse path ensembles as probability
distributions $P_{\text{fwd}}$ and $P_{\text{rev}}$ on trajectory
space. The Bhattacharyya coefficient between them satisfies
(by Jensen's inequality applied to Eq.~\eqref{eq:crooks_average}):
\begin{equation}
\BC(P_{\text{fwd}}, P_{\text{rev}})
\;\geq\; \exp\!\bigl(-\Sigma_{\text{Sch}}/2\bigr),
\label{eq:BC_crooks}
\end{equation}
where $\Sigma_{\text{Sch}} = \langle \Sigma_{\text{traj}} \rangle$
is the mean entropy production, equal to the KL divergence
$D_{\mathrm{KL}}(P_{\text{fwd}} \| P_{\text{rev}})$ on path space.

Since the Bhattacharyya coefficient is the classical fidelity
(Layer~3, Eq.~\eqref{eq:BC}), the classical $\tauasym$ bound
applies:
\begin{equation}
\tauasym_{\text{path}} \;\leq\;
1 - \exp\!\bigl(-\Sigma_{\text{Sch}}/2\bigr).
\label{eq:tau_path}
\end{equation}

\subsection{Esposito--Van den Broeck decomposition}

The total entropy production decomposes
as~\cite{EspositoVdB2010}:
\begin{equation}
\Sigma_{\text{total}}
= \Sigma_{\text{ad}} + \Sigma_{\text{na}},
\label{eq:EVdB}
\end{equation}
where the adiabatic contribution
\begin{equation}
\Sigma_{\text{ad}}
= D_{\mathrm{KL}}\bigl(p(t) \,\big\|\, p_{\text{eq}}(\lambda(t))\bigr)
\label{eq:Sigma_ad}
\end{equation}
measures the distance between the instantaneous distribution $p(t)$
and the equilibrium distribution $p_{\text{eq}}(\lambda(t))$
for the current protocol value. This is precisely a KL divergence
and is therefore the JRSWW-compatible piece: it equals the QRE drop
for the thermalization channel. The nonadiabatic contribution
$\Sigma_{\text{na}}$ captures the lag between the system's response
and the driving protocol.

\subsection{Application: brain entropy production}

Sanz Perl et~al.~\cite{SanzPerl2021} measured entropy production in
human brain dynamics using the KL divergence between forward and
time-reversed neural trajectories:
\begin{equation}
\Sigma_{\text{brain}} = D_{\mathrm{KL}}(P_{\text{fwd}}\|P_{\text{rev}}).
\label{eq:brain_sigma}
\end{equation}
They found $\Sigma_{\text{brain}}$ monotonically correlated with
consciousness level: highest in wakefulness, reduced under
anesthesia, lowest in unresponsive states. The $\tauasym$ bound
\eqref{eq:tau_path} applies to these measurements, giving an upper
limit on the retrodictability of brain states.

\subsection{IIT's $\Phi$ is not $\Sigma$}

It is tempting to identify integrated information $\Phi$
(from Integrated Information Theory~\cite{Tononi2008}) with
entropy production $\Sigma$. This identification is wrong.
Explicit counterexamples exist in both directions:
\begin{itemize}
\item High $\Phi$, low $\Sigma$: a system with strong integration
but operating near equilibrium (e.g., a balanced neural circuit
with slow dynamics).
\item High $\Sigma$, low $\Phi$: a highly dissipative but modular
system (e.g., two independent heat engines running in parallel).
\end{itemize}
$\Phi$ measures the irreducibility of a system's causal structure;
$\Sigma$ measures thermodynamic irreversibility. These are distinct
quantities that can vary independently.

%=======================================================================
\section{Summary}
\label{sec:summary}
%=======================================================================

Table~\ref{tab:summary} summarizes the verification across all four
layers.

\begin{table*}[t]
\caption{Mathematical verification of the $\tauasym$ bound across
four layers. Each layer is assigned a status label reflecting the
current state of proof.}
\label{tab:summary}
\begin{ruledtabular}
\begin{tabular}{lllll}
\textbf{Layer} & \textbf{Domain} & \textbf{Status}
& \textbf{Key equation} & \textbf{What it proves} \\
\hline
1 & Quantum channels
& \textsc{Theorem}
& $F \geq \exp(-\Sigma/2)$
& $\tauasym$ bound holds universally \\
2 & Gravitational fields
& \textsc{Motivated Conjecture}
& $\Sigma_{\text{grav}} = -\ln(-g_{00})$
& 3 routes converge \\
3 & Classical systems
& \textsc{Theorem}
& $\tauasym_{\text{cl}} \leq 1 - \exp(-\Delta D_{\mathrm{KL}}/2)$
& Special case of Layer~1 \\
4 & Stochastic thermo
& \textsc{Bridge Exists}
& $\tauasym_{\text{path}} \leq 1 - \exp(-\Sigma_{\text{Sch}}/2)$
& Via Crooks FT \\
\end{tabular}
\end{ruledtabular}
\end{table*}

The $\tauasym$ parameter has a secure mathematical foundation
(Layer~1), a well-motivated gravitational extension (Layer~2),
a rigorous classical limit (Layer~3), and a sound connection to
stochastic thermodynamics (Layer~4).

What remains open:
\begin{enumerate}
\item Layer~2 requires either a first-principles CPTP map for
gravity or a proof that the quantum hypothesis testing route
reproduces $\Sigma_{\text{grav}} = -\ln(-g_{00})$.
\item Layer~4 would benefit from an explicit derivation of
the nonadiabatic contribution's effect on the $\tauasym$ bound.
\end{enumerate}

%=======================================================================

\begin{thebibliography}{30}

\bibitem{JRSWW2018}
M.~Junge, R.~Renner, D.~Sutter, M.~M. Wilde, and A.~Winter,
Ann.\ Henri Poincar\'e \textbf{19}, 2955 (2018).

\bibitem{Uhlmann1976}
A.~Uhlmann,
Rep.\ Math.\ Phys.\ \textbf{9}, 273 (1976).

\bibitem{FawziRenner2015}
O.~Fawzi and R.~Renner,
Commun.\ Math.\ Phys.\ \textbf{340}, 575 (2015).

\bibitem{Petz1988}
D.~Petz,
Q.\ J.\ Math.\ \textbf{39}, 97 (1988).

\bibitem{HJPW2004}
P.~Hayden, R.~Jozsa, D.~Petz, and A.~Winter,
Commun.\ Math.\ Phys.\ \textbf{246}, 359 (2004).

\bibitem{Paper2}
S.-K. Huang,
``Gravity as quantum retrodiction cost,''
companion paper (2026).

\bibitem{DorauMuch2025}
R.~Dorau and A.~Much,
Phys.\ Rev.\ Lett.\ \textbf{134}, 011601 (2025).

\bibitem{Herrera2020}
L.~Herrera,
Entropy \textbf{22}, 1110 (2020).

\bibitem{AhmadiFuentes2014}
M.~Ahmadi, D.~E. Bruschi, and I.~Fuentes,
Phys.\ Rev.\ D \textbf{89}, 065028 (2014).

\bibitem{Wilde2017}
M.~M. Wilde, M.~Tomamichel, and M.~Berta,
IEEE Trans.\ Inf.\ Theory \textbf{63}, 1792 (2017).

\bibitem{GoldenThompson}
S.~Golden,
J.\ Math.\ Phys.\ \textbf{6}, 1188 (1965);
C.~J. Thompson,
J.\ Math.\ Phys.\ \textbf{6}, 1812 (1965).

\bibitem{HayashiYamasaki2025}
M.~Hayashi and H.~Yamasaki,
Nat.\ Phys.\ \textbf{21}, 1112 (2025).

\bibitem{Parzygnat2023}
A.~J. Parzygnat and F.~Buscemi,
Quantum \textbf{7}, 1013 (2023).

\bibitem{Umegaki1962}
H.~Umegaki,
Kodai Math.\ Sem.\ Rep.\ \textbf{14}, 59 (1962).

\bibitem{Schnakenberg1976}
J.~Schnakenberg,
Rev.\ Mod.\ Phys.\ \textbf{48}, 571 (1976).

\bibitem{Crooks1999}
G.~E. Crooks,
Phys.\ Rev.\ E \textbf{60}, 2721 (1999).

\bibitem{EspositoVdB2010}
M.~Esposito and C.~Van den Broeck,
Phys.\ Rev.\ Lett.\ \textbf{104}, 090601 (2010).

\bibitem{SanzPerl2021}
Y.~Sanz Perl, H.~Bocaccio, I.~P\'erez-Ipi\~na,
F.~Zamberlan, S.~Pallavicini, M.~Tagliazucchi,
K.~Kringelbach, G.~Deco, and E.~Tagliazucchi,
Phys.\ Rev.\ E \textbf{104}, 014120 (2021).

\bibitem{Tononi2008}
G.~Tononi,
Biol.\ Bull.\ \textbf{215}, 216 (2008).

\bibitem{Landi2021}
G.~T. Landi and M.~Paternostro,
Rev.\ Mod.\ Phys.\ \textbf{93}, 035008 (2021).

\bibitem{Buscemi2024}
F.~Bai, F.~Buscemi, and V.~Scarani,
arXiv:2412.12489 (2024).

\bibitem{Kwon2019}
H.~Kwon and M.~S. Kim,
Phys.\ Rev.\ X \textbf{9}, 031029 (2019).

\bibitem{Paper1}
S.-K. Huang,
``The arrow of time from Petz recovery,''
companion paper (2026).

\end{thebibliography}

\end{document}
